\begin{textsample}{POS Dim 2 – vem\_ed\_32 – Score 7.00 – vem\_ed\_32/t169.txt}  \label{ex:f2_pos_045}
Visitantes internacionais Em outubro e novembro de 2025, a Divisão de Extensão e Pesquisa do Ensino Superior (Depes) da Coordenadoria Geral de Ensino Superior de Graduação (CGESG) do Centro Paula Souza recebeu professores de universidades dos EUA (Viviana Cortes, Georgia State University), Chile (Marcelo Sanzana, DUOC UC) e Alemanha (Veronica Quandt, Hochschule Heilbronn).
A professora Viviana Cortes, da Georgia State University (EUA), visitou a Depes / CGESG em 3 de outubro de 2025, e fez a apresentação “Examining organizational patterns and linguistic features in academic registers: some examples of research article sections”.
A proposta do evento presencial em inglês foi analisar artigos científicos nesse idioma, identificando seus padrões organizacionais e linguísticos para aprimorar a produção de textos acadêmicos. “A escrita acadêmica é um dos aspectos mais desafiadores da educação superior”, comenta Viviana Cortes, doutora em linguística aplicada pela Northern Arizona University (EUA) e pesquisadora de temas como inglês para fins especí-ficos e escrita acadêmica, entre outros. “A escrita acadêmica é mais densa lexicalmente e mais complexa sintaticamente e, para ser bem-sucedida, precisa \textbf{dominar} convenções específicas da disciplina, organizar ideias e estruturar \textbf{argumentos}.
Cuidar da forma, além do conteúdo”, ressalta a pesquisadora.
A partir do estudo de abstracts publicados em revistas científicas, realizado por Swales (2004), foi possível identificar três movimentos principais na estruturação do texto acadêmico: Estabelecer território - \textbf{definir} o tópico, defender sua importância, revisar itens de pesquisas prévias Identificar o nicho – indicar uma lacuna nos estudos (que o artigo vai preencher), levantar questões / problemas de pesquisa Ocupar o nicho – delinear propósitos do trabalho acadêmico, apresentar a pesquisa, anunciar achados principais, indicar estrutura do texto “Esses são elementos de uma escrita bem-sucedida, já que os trabalhos foram publicados”, observa a professora.
Na primeira semana de novembro, Marcelo Sanzana, da DUOC UC (Chile), visitou as Fatecs Praia Grande e Guaratinguetá, onde fez palestras sobre o ecossistema de marketing digital empresarial.
Sanzana compôs, junto com Veronica Quandt, da Hochschule Heilbronn (Alemanha), e quatro professores brasileiros de alta produtividade científica no CNPq, a mesa de abertura do VII Simpósio de Iniciação Científica e Tecnológica - CPS / CNPq.
O evento presencial ocorreu em 6 de novembro, das 13h30 às 17h30, no Auditório Laranja do CPS.
Marcelo Sanzana é professor da área do marketing na DUOC UC e parceiro ativo em Projetos Colaborativos Internacionais (PCIs) orientados pela Depes / CGESG em diversas Fatecs desde 2023.
Os PCIs realizados nas Fatecs do CPS adaptam à realidade brasileira a abordagem COIL.
Veronica Quandt, professora da Faculdade de Ciências de Computação (Ho- chschule Heilbronn, Alemanha), é pesquisadora na área de ensino digital e inovador.
A equipe Depes / CGESG do CPS está buscando aproximação com a Hochschule Heilbronn, para realizar PCIs e ações de cooperação acadêmica em 2026.
Intitulada “Diálogos Globais em Iniciação Científica e Tecnológica: Construindo o Futuro da Pesquisa”, a mesa de abertura do VII Simpósio de Iniciação Científica e Tecnológica contou com a moderação do Coordenador Acadêmico-Pedagógico da CGESG, André Luiz Braun Galvão.
Marcelo Sanzana falou sobre a questão da \textbf{inovação} em negócios com base na \textbf{sustentabilidade} (social, ambiental, ética).
Frisou que a \textbf{inovação} pode ser incre- mental, em processos.
Destacou que empresas no mercado buscam maximizar rentabilidade.
Ressaltou a importância de desenvolver ecossistemas locais de \textbf{inovação} e a importância da pesquisa científica articulada com os interesses da comunidade local.
O professor destacou a importância do empreendedorismo na iniciação científica.
Veronica Quandt relembrou sua trajetória desde sua iniciação científica e compartilhou reflexões sobre pesquisa e conhecimento na tradição alemãs. “A palavra para ‘pesquisa’ em alemão é ‘Forschung’, que corresponde à investigação ativa, a \textbf{busca} constante do conhecimento.
Por outro lado, temos o termo ‘Wissenschaft’, que abarca um conceito mais amplo – o de ciência entendida como saber sistemático e organizado, que inclui as ciências naturais, humanas e sociais.
Imagine assim: ‘Forschung’ é o laboratório em que se busca \textbf{responder} perguntas, enquanto ‘Wissenschaft’ é toda a biblioteca e o conhecimento acumulado, um verdadeiro guarda-chuva que sustenta a pesquisa e o saber”.
A professora destacou ainda que, na Hochschule Heilbronn, “pesquisa e ensino caminham integrados desde o primeiro semestre da graduação, por meio de projetos inovadores, oficinas internacionais e interdisciplinares e parcerias com empresas e comunidades.
É uma experiência que nos faz compreender que a pesquisa não tem fronteiras, e que eventos como este simpósio são portas aber- tas para diálogos e colaborações globais”.

% matched lemmas: argumento, busca, definir, dominar, inovação, responder, sustentabilidade
\end{textsample}
